% sec-10: how the document was built, and what survives it.  Figure 20.
\section{Build Method and Reproducibility}

\subsection{The Eight-Stage Build}

The paper was built as one schedule of eight stages: five diagram-specification
stages, then a draft, a full and a final paper stage. Two rules shaped it. One
commit per file at the moment the file is written, so a reader can watch the
branch rather than wait for it. And no Python file anywhere, so the repository's
three lint-and-format jobs stay green; the diagrams stage emits a
machine-readable specification instead, and every tile it describes is drawn in
LaTeX by a vector glyph macro.

\begin{apptable}
\begin{tabularx}{\textwidth}{>{\raggedright\arraybackslash}p{1.3cm} >{\raggedright\arraybackslash}p{3.9cm} >{\raggedright\arraybackslash}X >{\raggedright\arraybackslash}p{1.7cm}}
\headrow \thead{Stage} & \thead{Output directory} & \thead{Figures or sections produced} & \thead{Commits}\\
\midrule
1 & \appfile{mermaid/} & Figures 1, 7, 12, 13, 19 & 6 \\
2 & \appfile{plantuml/} & Figures 6, 10, 15 & 4 \\
3 & \appfile{d2/} & Figures 2, 5, 8, 11, 16 & 6 \\
4 & \appfile{diagrams-python/} & Figures 4, 18, 20 & 4 \\
5 & \appfile{graphviz/} & Figures 3, 9, 14, 17 & 5 \\
6 & \appfile{draft-capital/} & Twelve sections, twenty sized slots & 10 or more \\
7 & \appfile{full-capital/} & Twenty figures drawn, twenty-one tables & 10 or more \\
8 & \appfile{final-capital/} & The senior-author pass & 10 or more \\
\end{tabularx}
\end{apptable}
\vspace{-0.65cm}
\tabcap{The eight stages, their output directory and their commit count. There
is no publication subdirectory at stage 8, by instruction, and no stage batches
its work to the end.}

\subsection{Why the Diagram Split Is Uneven}

Five mermaid, three plantuml, five d2, three diagrams, four graphviz. The split
follows purpose rather than quota, and the two five-counts fall where this paper
argues most often. Mermaid answers what happens next and how long, and a
capitalization plan is mostly a claim about sequence. D2 answers what contains
what and what is priced against what, and the rest of the plan is tables. The
per-figure reasoning is in the five directory READMEs and is not restated here.

\subsection{What Survives a Stop}

\begin{appfloat}[!tb]
\begin{appfig}[diagrams-python-type / custody clusters]
% The public cluster is drawn larger than the other three combined, because it
% holds more and because that proportion is the argument.  Its border is solid
% pablue1 while the other three are dashed: solid is permanent, dashed is
% conditional.
\dgnodew{dgtiled}{0}{0}{\glyphdoc}{Protocols, Phase 1 and 2}{v1}
\dgnode{dgtile}{2.7}{0}{\glyphflask}{IND package, as deposited}{v2}
\dgnodew{dgtilem}{5.4}{0}{\glyphcpu}{QSP stack and VVUQ suite}{v3}
\dgnode{dgtile}{0}{-2.2}{\glyphchart}{Twelve milestone artifacts}{v4}
\dgnodew{dgtiled}{2.7}{-2.2}{\glyphlock}{Replay bundles and hashes}{v5}
\begin{scope}[on background layer]
\node[dgcluster,draw=pablue1,solid,line width=0.9pt,
      fit=(v1)(v1l)(v2)(v2l)(v3)(v3l)(v4)(v4l)(v5)(v5l),inner sep=7pt] (cpub) {};
\end{scope}
\node[dgctitle,anchor=south west] at ([yshift=1.2mm]cpub.north west)
  {Public custodian, Zenodo and GitHub, permanent};
\node[d2cellg,minimum width=24mm,text width=21mm,anchor=north east]
  at ([xshift=-1.5mm,yshift=-1.5mm]cpub.north east) {no retention limit};
\dgnodeg{dgtileg}{8.70}{0}{\glyphmon}{Monitoring reports}{w1}
\dgnodeg{dgtileg}{11.40}{0}{\glyphdoc}{Forms 3454 and 3455}{w2}
\dgnodeg{dgtileg}{8.70}{-2.2}{\glyphlink}{FDA correspondence}{w3}
\begin{scope}[on background layer]
\node[dgcluster2,fit=(w1)(w1l)(w2)(w2l)(w3)(w3l),inner sep=7pt] (cspn) {};
\end{scope}
\node[dgctitle2,anchor=south west] at ([yshift=1.2mm]cspn.north west)
  {Sponsor file, ChemicalQDevice};
\node[d2cellg,minimum width=24mm,text width=21mm,anchor=north east]
  at ([xshift=-1.5mm,yshift=-1.5mm]cspn.north east) {two years, 312.57};
\dgnodeg{dgtilek}{0}{-5.0}{\glyphdoc}{Source documents}{x1}
\dgnodeg{dgtilek}{2.7}{-5.0}{\glyphhand}{Consent forms}{x2}
\dgnodeg{dgtilek}{5.4}{-5.0}{\glyphdb}{Medical records}{x3}
\begin{scope}[on background layer]
\node[dgcluster2,fit=(x1)(x1l)(x2)(x2l)(x3)(x3l),inner sep=7pt] (csit) {};
\end{scope}
\node[dgctitle2,anchor=south west] at ([yshift=1.2mm]csit.north west)
  {Site file, UC San Diego};
\node[d2cellg,minimum width=24mm,text width=21mm,anchor=north east]
  at ([xshift=-1.5mm,yshift=-1.5mm]csit.north east) {per policy, 312.62};
\dgnodeg{dgtileg}{8.70}{-5.0}{\glyphbank}{IND and amendments}{y1}
\dgnodeg{dgtileg}{11.40}{-5.0}{\glyphsignal}{Expedited safety reports}{y2}
\begin{scope}[on background layer]
\node[dgcluster2,fit=(y1)(y1l)(y2)(y2l),inner sep=7pt] (cfda) {};
\end{scope}
\node[dgctitle2,anchor=south west] at ([yshift=1.2mm]cfda.north west)
  {Agency file, FDA};
\dgnodew{dgtiled}{14.10}{-2.5}{\glyphteam}{Third party}{z1}
\draw[dgedgeb,line width=1pt] (cpub.east) to[bend left=16]
  node[font=\tiny\sffamily,fill=protowhite,inner sep=1.5pt,pos=0.55]
  {reproducible without any custodian} (z1.north west);
\draw[dgedged] (cspn.south) -- node[font=\tiny\sffamily,fill=protowhite,inner sep=1.5pt,pos=0.5]
  {on request only} (cfda.north);
\draw[dgedged] (csit.east) -- node[font=\tiny\sffamily,fill=protowhite,inner sep=1.5pt,pos=0.5]
  {never leaves the site} (cfda.west);
\draw[dgedged] (csit.south east) -- (9.60,-6.90);
\pxmark{9.90}{-6.90}
\node[font=\tiny\sffamily,text=protogray,anchor=west] at (10.30,-6.90)
  {no participant-level path to a third party};
\node[d2cellk,text width=24mm,minimum height=9mm] at (14.10,-5.20)
  {five of seven work products reproducible};
\node[font=\tiny\itshape,text=protogray,anchor=north west,text width=132mm]
  at (0,-7.55) {The public cluster's border is solid and the other three are
  dashed: solid is permanent, dashed is conditional on somebody continuing to
  hold the file. The blocked edge is drawn to a third of its length and
  terminated, so it reads as a path that does not complete.};
\end{appfig}
\vspace{-0.65cm}
\figcaption{Four custodians and what each still holds if the programme stops.\\
Everything a third party needs to reproduce the computation sits\\
with the first custodian, which needs no one to stay in business.}
\end{appfloat}

Five of seven work products are fully reproducible from the public custodian
alone: the ten-arm simulation, the 1000-patient digital twin, the 55-test
verification suite with its credibility score, the replay of any advisory
decision, and this paper with all twenty figures. The interlock bench
verification is partly reproducible, since the report and the manifest are
public but the rig is not. Only participant-level clinical results are
unreachable, because the source documents never leave the site
\cite{ecfr2026part312records}.

That is the argument this figure exists to make, and it is made once. If the
company dissolves at any milestone, the computational work does not become
unavailable, because it was never in the company's custody in the first place
\cite{repo2026}.

\subsection{Verification of This Document}

Every source set in this build was compiled with pdfLaTeX twice and BibTeX
before its commit, at zero errors, zero overfull boxes, zero undefined
citations, and zero undefined references. No raster image is used anywhere in
this paper: every figure is pure TikZ on an eight-token palette with no black
fill, and the palette audit is a search for the deleted near-black fill token,
which returns nothing.

The figure-to-caption distance is identical at 6.06 pt for all twenty figures
and all twenty-one tables. The audit that proves it is arithmetic rather than
visual: the count of \texttt{\textbackslash vspace\{-0.65cm\}} in the section
sources equals the count of \texttt{\textbackslash figcaption} plus the count of
\texttt{\textbackslash tabcap}, and the style file closes both carriers with a
single rigid skip that no section can override.

Three defects were found and fixed during the build rather than carried forward,
and they are recorded because a verification section that reports only successes
is not a record. The three-cell cover ledger exceeded the measure by 14.34 pt,
because three cells at 0.288 of the text width plus two gaps plus six inner
separations do not fit inside it. A directory path set in a monospace face
overflowed a 3.3 cm table column by 11.23 pt in three rows, because Courier
offers no break point inside a path. And a TikZ node width written as
\texttt{\{\textbackslash w*0.415\} cm} is parsed by the mathematics engine as
one expression and fails; the unit belongs inside the braces.
